\section{Strongly minimal geometry and the Baldwin--Lachlan refinement}
\label{sec:p03-n-a-006-02}

The refinement of uncountable categoricity by Baldwin and Lachlan uses
strongly minimal sets rather than identifying cardinality with an ordinal
rank. Their 1971 article treats the strongly minimal approach to categoricity
and then proves the countable-model dichotomy
\cite{BaldwinLachlan1971}. The geometric results below have explicit local
proofs. Their hypotheses do not include the assertion that every
uncountably categorical theory has already been coordinatized by such a
set: that assertion is a further theorem, not a definition of dimension.

\begin{hbconvention}{def:p03-bl-ambient}{Parameters and algebraicity}
Fix a complete first-order theory \(T\), a model \(\mathcal U\), and a finite
tuple \(\bar c\) in \(\mathcal U\). Let
\(D=\delta(\mathcal U,\bar c)\) be infinite and strongly minimal: in every
elementary extension of \((\mathcal U,\bar c)\), every subset of the
interpretation of \(D\) definable with arbitrary parameters is finite or
cofinite. All uses of \(\operatorname{acl}\) refer to the ambient structure
and its elementary extensions, not just to the language induced on \(D\).
For \(B\subseteq\mathcal U\), an element is algebraic over \(B\) if a
formula with a finite tuple of parameters from \(B\) defines a finite set
containing it. For \(X\subseteq D\), put
\[
 \operatorname{cl}_{\bar c}(X)
   \defeq D\cap\operatorname{acl}(\bar c\cup X).
\]
The notation \(\bar c\cup X\) means the set of coordinates of \(\bar c\)
together with \(X\). A set \(I\subseteq D\) is independent over a parameter
set \(B\supseteq\bar c\) when
\(a\notin\operatorname{acl}(B\cup(I\setminus\{a\}))\) for every \(a\in I\).
A basis of \(D\) over \(\bar c\) is a maximal such independent subset.
\end{hbconvention}

\begin{hblemma}{thm:p03-bl-uniform}{proved=annex:p3.n-a-006.2.uniform}{Uniform finite/cofinite bounds}
For every formula \(\theta(x,\bar y,\bar c)\), where \(x\) ranges over \(D\),
there is a natural number \(q\) such that, in every elementary extension and
for every \(\bar b\), either
\[
 \abs{\{x\in D:\theta(x,\bar b,\bar c)\}}\leq q
 \quad\text{or}\quad
 \abs{\{x\in D:\neg\theta(x,\bar b,\bar c)\}}\leq q.
\]
In particular, within this family infinitude, cofinitude, and having more
than \(q\) solutions are equivalent and are expressible by a formula in
\(\bar y\).
\end{hblemma}

The elementary-extension clause in strong minimality is essential here.
Compactness must exclude a parameter instance with infinitely many
solutions and infinitely many nonsolutions, not just instances visible in
one small model.

\begin{hbtheorem}{thm:p03-bl-pregeometry}{proved=annex:p3.n-a-006.2.pregeometry}{Algebraic closure gives a pregeometry}
The operation \(\operatorname{cl}_{\bar c}\) is extensive, monotone,
idempotent, and of finite character. For \(a,b\in D\) and \(X\subseteq D\),
it satisfies exchange:
\[
 a\in\operatorname{cl}_{\bar c}(X\cup\{b\})
       \setminus\operatorname{cl}_{\bar c}(X)
 \quad\Longrightarrow\quad
 b\in\operatorname{cl}_{\bar c}(X\cup\{a\}).
\]
The conclusion remains valid with any additional fixed parameter set in
the algebraic-closure base. No countability hypothesis is needed for this
theorem.
\end{hbtheorem}

The proof is a finite-fibre argument. Failure of exchange would make a
relation have cofinite fibres in one direction and uniformly bounded finite
fibres in the other. Intersecting finitely many cofinite fibres then
contradicts the finite bound. This is a proof of the difficult closure
axiom, not an assumption that the set is already a vector space.

\begin{hbtheorem}{thm:p03-bl-dimension}{proved=annex:p3.n-a-006.2.dimension}{Bases and dimension}
Assume the Axiom of Choice in the ambient set theory. Every independent
subset of \(D\) extends to a basis, each basis \(I\) spans
\(D=\operatorname{cl}_{\bar c}(I)\), and all bases have the same cardinality.
Write this cardinal as \(\dim_{\bar c}(D)\). More generally, an independent
set in the closure of a finite set of \(r\) elements has at most \(r\)
elements. If the language is countable and \(\bar c\) is finite, then
\[
 \abs{D}\leq\max(\aleph_0,\dim_{\bar c}(D)).
\]
Consequently an uncountable \(D\) has
\(\dim_{\bar c}(D)=\abs{D}\).
\end{hbtheorem}

\begin{hblemma}{thm:p03-bl-independent-types}{proved=annex:p3.n-a-006.2.independent-types}{The type of an independent tuple}
For every parameter set \(B\supseteq\bar c\), all elements of
\(D\setminus\operatorname{acl}(B)\) have the same complete type over \(B\).
Any two ordered independent tuples in \(D\) of the same finite length over
\(B\) have the same complete type over \(B\). Hence any bijection between
independent subsets over \(B\), together with the identity on \(B\), is an
elementary map on its domain.
\end{hblemma}

\subsection{The whole-universe case and its exact limit}

\begin{hbdefinition}{def:p03-bl-strongly-minimal-theory}{A strongly minimal theory}
A complete theory \(T\) with infinite models is strongly minimal when
\(x=x\) is strongly minimal in its models. Thus its entire universe, not
merely one definable subset, is a strongly minimal set. In this subsection
the base tuple is empty and every model \(M\) has a basis for
\(\operatorname{acl}\) on its whole universe.
\end{hbdefinition}

\begin{hbtheorem}{thm:p03-bl-whole-models}{proved=annex:p3.n-a-006.2.whole-models}{Classification and prime-over-basis in the whole-universe case}
Assume ambient choice, and let \(T\) be strongly minimal. Two models are isomorphic if and only if
their bases have the same cardinality; any bijection between bases extends
to an isomorphism. Every model \(M\) is prime and minimal over a basis
\(I\): any elementary map of \(I\) into a model of \(T\) extends to an
elementary embedding of \(M\), and an elementary submodel of \(M\)
containing \(I\) is \(M\) itself.

For a countable language, \(T\) is categorical in every uncountable
cardinal, and its countable-model spectrum is either one or \(\aleph_0\).
More precisely, if a finite-dimensional model exists, there is a least
finite dimension \(n_0\), and the countable models have exactly the
dimensions \(n_0,n_0+1,\ldots,\aleph_0\). If no finite-dimensional model
exists, every countable model has dimension \(\aleph_0\).
\end{hbtheorem}

\begin{proofboundary}{def:p03-bl-coordinatization-boundary}{Geometry is not yet general coordinatization}
The preceding construction is justified by
\(M=\operatorname{acl}(I)\) for the \emph{whole universe}. It does not
prove that, for an arbitrary uncountably categorical theory, a strongly
minimal \(D\) exists over parameters in a prime model or that every model
is prime and minimal over a basis of \(D\). Nor does prime-over-a-basis
imply that all elements are algebraic over that basis in this more general
setting. The existence and model-control assertions belong to the
Baldwin--Lachlan refinement and require their own argument
\cite{BaldwinLachlan1971}. The original article's general countable-model
argument and its fourth-section example are not replaced by the
whole-universe proof here.
\end{proofboundary}

\begin{hbexample}{def:p03-bl-extra-sort}{A strongly minimal subset that does not control the model}
Take the language with one unary predicate \(P\) and the complete theory
that both \(P\) and its complement are infinite pure sets. Its quantifier
elimination follows by extending finite partial bijections separately in
the two cells. The set \(P\) is strongly minimal: with finitely many
parameters, a definable subset of \(P\) is a finite or cofinite union of
equality cells. Nevertheless the full model also has the independent
cardinal invariant \(\abs{M\setminus P}\). In an uncountable cardinal
\(\kappa\), the models with cell sizes \((\kappa,\aleph_0)\) and
\((\kappa,\kappa)\) are not isomorphic, even though their \(P\)-dimensions
agree. The mere presence of a strongly minimal definable set therefore
does not supply the model-control step.
\end{hbexample}

\subsection{Fixed-characteristic algebraically closed fields}

\begin{hbtheorem}{thm:p03-bl-acf-geometry}{proved=annex:p3.n-a-006.2.acf}{Field algebraicity and the dimension illustration}
Fix \(p=0\) or a prime, and work in the ring language with
\(\mathsf{ACF}_p\). Using the quantifier-elimination theorem
\ref{thm:algebraically-closed-fields}, this is a complete strongly minimal theory.
If \(K\models\mathsf{ACF}_p\) and \(B\subseteq K\), then
\(\operatorname{acl}(B)\) is the set of elements of \(K\) algebraic over
the field generated by \(B\) and the prime field. A pregeometry basis is
therefore a transcendence basis. There is exactly one countable model for
each transcendence degree \(0,1,2,\ldots,\aleph_0\), and exactly one model
in each uncountable cardinal.
\end{hbtheorem}

For an infinite one-variable definable set in a strongly minimal universe,
Morley rank is \(1\) and Morley degree is \(1\): every definable proper
split has a finite side, and two cofinite subsets cannot be disjoint.
These invariants stay fixed while the dimension of a countable
algebraically closed field ranges through \(0,1,2,\ldots,\aleph_0\).
Thus neither rank nor degree is the dimension of the field's basis.
Characteristic must be fixed; fields of different characteristics are
not elementarily equivalent. This field-algebraicity illustration uses
the quantifier-elimination and algebraic-closure dependencies stated
above; it is distinct from the general Baldwin--Lachlan coordinatization
argument.

Morley's original categoricity argument and the strongly minimal
refinement are different source presentations. The former's rank and
saturation/transfer machinery is not renamed as a pregeometry proof, and
the local proofs above do not establish the refinement's general
existence theorem \cite{Morley1965,BaldwinLachlan1971}.
